\chapter{Expected Results}


This chapter outlines the anticipated findings and success criteria for evaluating the proposed intelligent caching framework. As this is a proposal, the values presented are educated estimates based on the system architecture and preliminary analysis. Actual measurements will be collected during the implementation phase and reported in the final thesis.

\section{Expected Caching Performance}

\subsection{Cache Hit Rates}

Based on the feature caching design, we expect the cache hit rates shown in Table~\ref{tab:expected_hit_rates}.

\begin{table}[ht]
\centering
\begin{tabular}{|l|c|l|}
\hline
\textbf{Scenario} & \textbf{Feature Cache Hit Rate} & \textbf{Explanation} \\
\hline
First Run (Cold Start) & 0\% & Features computed and cached \\
\hline
Identical Re-run & 100\% & Fingerprint matches \\
\hline
Changed Preprocessing & 0\% & Fingerprint changes, recompute \\
\hline
Changed Feature Selection Only & 100\% & Selection applied at query time \\
\hline
Add 20 New Subjects & 86\% & 128/148 subjects already cached \textcolor{red}{HIGHLIGHT FEATURES PERSISTENT !!! maybe delete here}  \\
\hline
\textbf{Typical Research Workflow} & \textbf{$>$90\%} & Model/hyperparam changes reuse cache \\
\hline
\end{tabular}
\caption{Expected feature cache hit rates across different experimental scenarios.}
\label{tab:expected_hit_rates}
\end{table}
\textcolor{red}{Also change 90 percent if 20 new subjects will be deleted }
The ``Typical Research Workflow'' row reflects that researchers often re-run experiments with minor parameter changes (e.g., different feature selection thresholds, different models), which should yield high cache hit rates for preprocessing and feature extraction.

\subsection{Measured Time Savings}

Table~\ref{tab:expected_time} presents the computation times and savings based on implementation benchmarks.

\begin{table}[ht]
\centering
\begin{tabular}{|l|c|c|c|}
\hline
\textbf{Operation} & \textbf{Cold Start} & \textbf{Cached} & \textbf{Speedup} \\
\hline
\multicolumn{4}{|c|}{\textit{Feature Extraction (Implemented -- Persistent Storage)}} \\
\hline
Preprocessing (per subject) & $\sim$5 sec & $\sim$5 sec & 1$\times$ (not cached) \\
\hline
Feature Extraction (per subject) & $\sim$25 sec & $\sim$0.1 sec & 250$\times$ \\
\hline
Full Pipeline (128 subjects) & $\sim$53 min & $\sim$14 sec & \textbf{227$\times$} \\
\hline
\multicolumn{4}{|c|}{\textit{LOSO Model Training (Proposed -- Intelligent Caching)}} \\
\hline
Single LOSO fold & $\sim$5 min & $\sim$1 sec & 300$\times$ \\
\hline
18-Config Experiment (2304 folds) & $\sim$16 hours & $\sim$4 min & \textbf{240$\times$} \\
\hline
\end{tabular}
\caption{Expected timing breakdown: Feature caching (implemented) vs. LOSO model caching (proposed). Feature extraction times are measured; LOSO model times are estimates based on preliminary benchmarks.}
\label{tab:expected_time}
\end{table}

\textbf{Key Insight:} The 227$\times$ speedup is achieved because the extraction of 149 EEG features per epoch (Stage 2) is cached once and reused across all 18 model configurations. This transforms a 16-hour experiment into a 4-minute experiment after the \textcolor{red}{complete} initial cold start. \textcolor{red}{check this section and think about keeoing feautre cache as time saved? better no}

\subsection{Measured Storage Requirements}

Table~\ref{tab:expected_storage} shows the actual storage overhead of the caching system using NPZ compression.

\begin{table}[ht]
\centering
\begin{tabular}{|l|c|c|}
\hline
\textbf{Cache Level} & \textbf{Per Subject} & \textbf{128 Subjects} \\
\hline
Stage 2 (149 Features, NPZ compressed) & $\sim$1.1 MB & $\sim$143 MB \\
\hline
\textbf{Total Cache} & $\sim$1.1 MB & \textbf{$\sim$143 MB} \\
\hline
\end{tabular}
\caption{Measured storage requirements for the feature caching system.}
\label{tab:expected_storage}
\end{table}



\begin{center}
\colorbox{yellow!20}{%
\begin{minipage}{0.9\textwidth}
\textbf{Working Draft Note:} 
Note outdated needs to be changed or deleted
\end{minipage}
}
\end{center}
\textcolor{red}{\textbf{Note:}} The current implementation caches only the extracted features (Stage 2). Preprocessed epochs (Stage 1) are computed on-the-fly during each run and not persisted to disk. This design prioritizes the feature extraction bottleneck, which accounts for $>$95\% of computation time. The compact NPZ format provides efficient compression using DEFLATE/zlib. For the 128-subject dataset with approximately 120,000 total epochs ($\sim$940 epochs per subject) and 195 features each, the uncompressed feature matrices would require approximately 183~MB (195 features $\times$ 120,000 epochs $\times$ 8 bytes = 187~MB). The compressed NPZ cache achieves $\sim$146~MB total, representing a compression ratio of approximately 1.25$\times$. This modest compression is typical for dense floating-point arrays, where the primary storage benefit comes from the efficient binary format rather than aggressive compression.

\section{Expected Classification Performance}


\begin{center}
\colorbox{yellow!20}{%
\begin{minipage}{0.9\textwidth}
\textbf{Working Draft Note:} 
should we keep this ? or weaken the benchmarks as this isnt our goal at all
\end{minipage}
}
\end{center}
While classification performance is secondary to the caching contribution, we expect results consistent with literature for EEG-based sleep staging.

\subsection{Overall Performance Expectations}

Based on the capstone project results and literature review, Table~\ref{tab:expected_classification} presents expected classification metrics.

\begin{table}[ht]
\centering
\begin{tabular}{|l|c|c|c|}
\hline
\textbf{Model} & \textbf{Accuracy} & \textbf{Macro F1} & \textbf{Cohen's Kappa} \\
\hline
XGBoost & 72--78\% & 58--68\% & 60--72\% \\
\hline
Random Forest & 68--75\% & 55--65\% & 55--68\% \\
\hline
\end{tabular}
\caption{Expected classification performance ranges based on LOSO cross-validation.}
\label{tab:expected_classification}
\end{table}

\subsection{Per-Class Performance Expectations}

Sleep stage classification exhibits known challenges with certain stages. Table~\ref{tab:expected_per_class} shows expected per-class F1 scores.

\begin{table}[ht]
\centering
\begin{tabular}{|l|c|l|}
\hline
\textbf{Stage} & \textbf{Expected F1} & \textbf{Notes} \\
\hline
Wake (W) & 75--85\% & Distinct EEG patterns \\
\hline
N1 (Light Sleep) & 20--40\% & Known difficulty---transitional stage \\
\hline
N2 (Sleep) & 75--85\% & Majority class, distinctive spindles \\
\hline
N3 (Deep Sleep) & 70--80\% & Strong delta waves \\
\hline
REM & 70--80\% & Distinctive mixed-frequency pattern \\
\hline
\end{tabular}
\caption{Expected per-class F1 scores with explanatory notes.}
\label{tab:expected_per_class}
\end{table}

\textbf{Note on N1 Classification:} The low expected F1 for N1 is consistent with literature and clinical practice. N1 is a brief transitional stage with EEG patterns similar to both Wake and N2, making it inherently difficult to classify. This is not a limitation of our approach but a known characteristic of the sleep staging problem.

\section{Expected Reproducibility Validation}

\subsection{Fingerprint Determinism}

We expect 100\% fingerprint determinism: the same configuration should always produce the same SHA-256 hash. This will be validated by running fingerprint generation multiple times with identical inputs.

\subsection{Result Consistency}

For identical fingerprints, we expect:
\begin{itemize}
\item Classification metrics identical within floating-point precision ($<$0.001\% variance)
\item Cached data byte-identical across loads (verified via checksum)
\item Model predictions identical for deterministic algorithms (XGBoost, RF with fixed seed)
\end{itemize}

\textbf{Note:} FNN models may exhibit minor variance ($\pm$0.5\%) due to GPU non-determinism, which will be documented as a known limitation.

\section{Success Criteria}




\begin{center}
\colorbox{yellow!20}{%
\begin{minipage}{0.9\textwidth}
\textbf{Working Draft Note:} 
Educated guesses so far:
\end{minipage}
}
\end{center}

The thesis will be considered successful if the following criteria are met:

\begin{enumerate}
\item \textbf{Caching Efficiency:}
  \begin{itemize}
  \item Cache hit rate = 100\% for repeated experiments with identical configurations
  \item Speedup factor $>$100$\times$ for cached features (achieved: 227$\times$)
  \item Storage overhead $<$200 MB for full 128-subject dataset (achieved: $\sim$143 MB)
    \item Speedup factor $>$10$\times$ for cached loso folds (128 subjects) (achieved: \textcolor{red}{placholder}$\times$)
  \end{itemize}

\item \textbf{Reproducibility:}
  \begin{itemize}
  \item 100\% fingerprint determinism
  \item Result variance $<$0.1\% for identical configurations
  \item Successful pipeline resumption after interruption
  \end{itemize}

\item \textbf{Classification Validity:}
  \begin{itemize}
  \item Overall accuracy $>$70\% (validates correct implementation)
  \item Macro F1 $>$55\% (handles class imbalance)
  \item Results consistent with published literature
  \end{itemize}
\end{enumerate}

